\section{Conclusion}

We introduced the \Core{}, a reusable, event‑driven substrate that
provides four essential self‑regulation capabilities: decoupled
communication, adaptive thresholds, chaos detection, and thermodynamic
cost accounting.  The core is implemented in Python, fully instrumented
with Prometheus, and released as open‑source software.  Benchmarks on a
commodity laptop (ASUS Strix GL704GW, Intel i7‑8750H, NVIDIA RTX 2070)
show that the bus introduces a per‑event overhead of only $6.9\,\mu$s,
while a single‑subscriber configuration sustains over 120\,k events/s.
The predictive warning mechanism of the ChaosDetector detected divergent
behaviour 78 time steps earlier than a classical threshold, and GPU‑based
computation reduced energy consumption by a factor of 24$\times$ compared
to CPU, demonstrating that the \Core{} enables both safe and
energy‑efficient autonomous operation.

The \Core{} is not a complete agent; it is the scaffold on which such
agents can be built.  Its design reflects the conviction that
autonomous knowledge systems must be \emph{constructively self‑aware}—
continuously monitoring, calibrating, and justifying their own
operation.